\section{Wnioski}
\begin{itemize}
\item Model liniowy dobrze przybliża obiekt w pobliżu punktu pracy. Przy dużych zmianach dopływu błędy poziomu i dynamiki stają się wyraźne.
\item DMC uwzględnia opóźnienie i zapewnia śledzenie zadania oraz tłumienie zakłóceń w badanych próbach. Przy zmianach zadania o $\pm8\ \mathrm{cm}$ odpowiedź nieliniowa jest bardziej oscylacyjna.
\item Jeden model lokalny nie opisuje równie dobrze całego zakresu pracy. Wyniki stanowią punkt odniesienia dla regulatorów rozmytych w dalszych częściach projektu.
\end{itemize}
\smallskip
{\small Źródła: treść zadania \texttt{main (50).pdf}, część 1; skrypt \texttt{run\_part1.m}.}
